% =====================================================================
% Section 4. Numerical results  (structure as Battisti & Boscheri 2025, Sec. 4)
% Source: src/spb6_physics.h (test_setup, test_state), tests/run_archive.sh, data/*/note.md
% Only the test descriptions and initial data are written; the results are to be added.
% =====================================================================
\section{Numerical results}
\label{sec:results}

All computations use the scheme of Sec.~\ref{sec:method} with the same settings: LSDIRK2, the material CFL condition
\eqref{eq:cfl} with $C=0.2$, minmod reconstruction, a local Rusanov speed on each face and the Newton tolerance
of Sec.~\ref{sec:fully}. Thermal and chemical relaxation are active in cells with $10^{-4}\le\alpha_1\le1-10^{-4}$. The
exceptions are stated with each test and collected in Table~\ref{tab:tests}. Where a phase is absent its volume
fraction is set to $10^{-8}$ or $10^{-6}$ as indicated; initial data in two dimensions are cell averages computed with
$4\times4$ sub-cells. Units are SI unless stated otherwise.

The relaxation levels are denoted as in Pelanti and Shyue (2014). $p$: instantaneous mechanical relaxation only.
$p$--$pT$: in addition, instantaneous thermal relaxation in interface cells. $p$--$pT$--$pTg$: in addition,
instantaneous mass transfer in interface cells where the $pT$-equilibrium temperature exceeds $T_{\rm sat}(p)$.
$p$--$pTg$: thermal and chemical relaxation together, only in interface cells where a phase temperature at
mechanical equilibrium exceeds $T_{\rm sat}(p)$. A finite rate replaces the corresponding instantaneous relaxation and
is marked $(f)$. Reference solutions are exact Riemann solutions of the relaxed limit models: the five-equation
model for $p$, the four-equation model for $p$--$pT$, and the $pT\to pTg$ wave fan and the homogeneous equilibrium
model for instantaneous mass transfer.

\begin{table}[t]
\centering\small
\caption{Test problems. $N$: number of cells; $C$: material CFL number; ``global'': Rusanov speed taken as the maximum
over the domain instead of per face.}
\label{tab:tests}
\resizebox{\textwidth}{!}{%
\begin{tabular}{llllll}
\toprule
Sec. & problem & EOS & $N$ & $t_{\rm end}$ & relaxation \\
\midrule
\ref{sec:res_vortex} & two-phase isentropic vortex & ideal gases & $32^2\ldots256^2$ & 5 & $p$ \\
\ref{sec:res_tg} & Taylor--Green vortex & ideal gases & $50^2$; $32^2\ldots256^2$ & 0.2 & $p$ \\
\ref{sec:res_interface} & water--air interface, $u=100$\,m/s & SG & 200 & 2.79\,ms & $p$ \\
 & water--air interface, $u=1$\,m/s & SG & 200; 1000 & 0.25\,s & $p$ \\
 & air density step, $u=1$\,m/s & ideal gas & 200; 1000 & 0.5\,s & -- \\
\ref{sec:res_single} & five single-phase Riemann problems & ideal gas; SG & 200; 1000 & Table~\ref{tab:single} & -- \\
\ref{sec:res_shock} & water--air, 1\,GPa & SG & 1000; 100 & 240\,$\mu$s & $p$ \\
 & water--air, 1\,GPa, $\rho_{\rm air}=50$ & SG & 200; 1000; 5000 & 241\,$\mu$s & $p$ \\
 & water--air, 1\,TPa & SG & 1000 & 8.3\,$\mu$s & $p$ \\
 & nitromethane & Cochran--Chan & 500 & 67\,$\mu$s & $p$ \\
 & detonation products--water & JWL / NASG & 500 & 0.9\,ms & $p$ \\
 & epoxy--spinel, $10^{10}$\,Pa ($C=0.05$) & SG & 500 & 80\,$\mu$s & $p$ \\
 & epoxy--spinel, $10^{10}$\,Pa & SG & 200; 1000 & 30\,$\mu$s & $p$ \\
 & epoxy--spinel, $2\times10^{11}$\,Pa (global) & SG & 500 & 34\,$\mu$s & $p$ \\
\ref{sec:res_cav} & water--air expansion tube & SG & 1000 & 1.85\,ms & $p$ \\
 & water cavitation tube, $\mp2$, $\mp500$\,m/s & SG & 5000 & 3.2, 0.58\,ms & four levels \\
 & water cavitation tube, $\mp2$\,m/s & NASG & 5000 & 3\,ms & five levels \\
\ref{sec:res_dodecane} & dodecane shock tube ($C=0.05$ with SG) & SG; NASG & 2000 & 473\,$\mu$s & $p$; $p$--$pT$--$pTg$; $p$--$pTg$ \\
\ref{sec:res_bartak} & Bart\'ak depressurisation & NASG & 1000 & 17.5\,ms & three cases \\
\ref{sec:res_he} & shock--helium bubble & ideal gases & $600\times160$ & 600\,$\mu$s & $p$ \\
\ref{sec:res_rm} & cavitating Richtmyer--Meshkov (global) & SG & $300\times100$ & 8.6\,ms & $p$ \\
\ref{sec:res_bubble} & vapour bubble compression (global) & SG & $200\times200$ & 1.2\,ms & $p$--$pT$--$pTg$ \\
\ref{sec:res_undex} & underwater explosion near a wall & JWL / NASG & $361\times200$ & 0.4\,ms & $p$ \\
\bottomrule
\end{tabular}}
\end{table}

% ---------------------------------------------------------------------
\subsection{Two-phase isentropic vortex}
\label{sec:res_vortex}
This test of Battisti and Boscheri measures the order of accuracy for a smooth two-phase flow at decreasing
Mach number. The domain is $[-5,5]^2$ with periodic boundaries; phase 1 is an ideal gas with $\gamma_1=1.35$ and
phase 2 an ideal gas with $\gamma_2=1.4$. With $r^2=x^2+y^2$, $\varepsilon=5$ and $\gamma=1.4$,
\begin{equation}
\begin{aligned}
&u=2-\frac{\varepsilon}{2\pi}e^{(1-r^2)/2}\,y,\qquad v=2+\frac{\varepsilon}{2\pi}e^{(1-r^2)/2}\,x,\qquad
T=p_0-\frac{(\gamma-1)\varepsilon^2}{8\gamma\pi^2}e^{1-r^2},\\
&\rho_1=\rho_2=\Big(\frac{T}{p_0}\Big)^{1/(\gamma-1)},\qquad p_1=p_2=\rho_1T,\qquad \alpha_2=\frac12+\frac14e^{(1-r^2)/2}.
\end{aligned}
\label{eq:vortex}
\end{equation}
The phases have equal pressure and density, the velocity field is divergence-free and $\alpha_2$ and the entropy are
constant along streamlines, so the exact solution is the translation of \eqref{eq:vortex} with velocity $(2,2)$; at
$t=5$ it coincides with the initial data. The background pressure is $p_0=1,10,10^2,10^3,10^4$, which gives
$M=2\sqrt2/\sqrt{\gamma p_0}$ from $2.4$ to $0.024$, and the grids have $N^2$ cells with $N=32,64,128,256$. Two initial
data are used. In the first the cell averages of \eqref{eq:vortex} are taken. In the second, well-prepared, the
velocity is the central difference of the stream function $\psi=\frac{\varepsilon}{2\pi}e^{(1-r^2)/2}$ at the cell
centres, $u=2+\delta_y\psi$, $v=2-\delta_x\psi$, which is discretely divergence-free; the other variables are cell
averages. In both cases the errors are measured in $L_2$ against the exact cell averages. As for every convergence
test in this section, the slopes are the unlimited central differences.

% ---------------------------------------------------------------------
\subsection{Low-Mach Taylor--Green vortex}
\label{sec:res_tg}
This test of Battisti and Boscheri checks the low-Mach limit of Sec.~\ref{sec:ap}. On $[0,2\pi]^2$ with periodic
boundaries both phases are ideal gases with $\gamma=1.4$, and
\begin{equation}
\rho_1=\rho_2=1,\qquad u=\sin x\cos y,\qquad v=-\cos x\sin y,\qquad p_1=p_2=p_0+\tfrac14(\cos2x+\cos2y),
\label{eq:tg}
\end{equation}
a stationary solution of the incompressible Euler equations. The volume fraction is $\alpha_2=\tfrac12$ (single-phase
case) or $\alpha_2=\tfrac12+\tfrac1{10}(\cos2x+\cos2y)$ (two-phase case). The background pressure
$p_0=10^2,10^4,\dots,10^{12}$ gives $M=1/\sqrt{1.4\,p_0}$ from $8.45\times10^{-2}$ to $8.45\times10^{-7}$. On a $50\times50$ grid the
solution is computed to $t=0.2$ with the fixed time step $\Delta t=2.5\times10^{-3}$, independent of $M$; the order of
accuracy is measured on $N=32,\dots,256$ with $\Delta t=0.08/N$. The errors are measured against \eqref{eq:tg}, together
with the discrete divergence of the face velocities.

% ---------------------------------------------------------------------
\subsection{Interface advection}
\label{sec:res_interface}
A water--air interface is advected in $[0,1]$ with transmissive boundaries (Saurel et al.\ 2009, Sec.~4.1). Water is a
stiffened gas with $\gamma=4.4$, $\pi=6\times10^8$\,Pa, air an ideal gas with $\gamma=1.4$. Initially $u=100$\,m/s,
$p=10^5$\,Pa, $\rho_{\rm water}=1000$, $\rho_{\rm air}=10$, and $\alpha_{\rm water}=1-10^{-8}$ for $x<0.5$, $10^{-8}$ otherwise. The
computation uses 200 cells and ends at $t=2.79$\,ms, when the interface is at $x=0.779$. Two low-Mach variants follow
Schropff et al.\ (2026): the same fluids with $u=1$\,m/s, $\rho_{\rm air}=0.01$, $\alpha_{\rm water}=1-10^{-7}$ for $x<0.25$ and $10^{-7}$
otherwise, $t_{\rm end}=0.25$\,s ($M\approx3\times10^{-3}$); and a single gas, air, with a density step $\rho=1000$ for $x<0.25$
and $\rho=0.01$ otherwise, $u=1$\,m/s, $p=10^5$\,Pa, $t_{\rm end}=0.5$\,s, so that $M$ ranges from $2.7\times10^{-4}$ to
$8.5\times10^{-2}$. Both use 200 and 1000 cells. In all three cases uniform velocity and pressure must be preserved
(Proposition~\ref{prop:wb}).

% ---------------------------------------------------------------------
\subsection{Single-phase Riemann problems}
\label{sec:res_single}
The Riemann problems of Table~\ref{tab:single} (Schropff et al.\ 2026) test the scheme for a single fluid over a wide
range of Mach numbers. A single fluid is represented by two identical phases in the same state, with
$\alpha_1=1-10^{-8}$. Air is an ideal gas with $\gamma=1.4$ and water the stiffened gas of Sec.~\ref{sec:res_interface}. The
domain is $[0,1]$ with transmissive boundaries, the initial discontinuity is at $x=0.5$, and 200 and 1000 cells are
used. The reference is the exact solution of the Euler equations.
\begin{table}[t]
\centering\small
\caption{Single-phase Riemann problems (Schropff et al.\ 2026): left and right states and final time.}
\label{tab:single}
\begin{tabular}{lllll}
\toprule
problem & fluid & $(\rho,u,p)_L$ & $(\rho,u,p)_R$ & $t_{\rm end}$ \\
\midrule
low-Mach shock tube & water & $(1000,\ 0,\ 10^8)$ & $(1000,\ 15,\ 0.98\times10^8)$ & $10^{-4}$ \\
liquid--liquid shock tube & water & $(1000,\ 0,\ 10^6)$ & $(1000,\ 0,\ 10^5)$ & $1.5\times10^{-4}$ \\
Lax shock tube & air & $(0.445,\ 0.698,\ 3.528)$ & $(0.5,\ 0,\ 0.571)$ & 0.14 \\
strong shock tube & air & $(1,\ 0,\ 10^{10})$ & $(0.125,\ 0,\ 0.1)$ & $2.5\times10^{-6}$ \\
two rarefactions & air & $(1,\ -1,\ 0.4)$ & $(1,\ 1,\ 0.4)$ & 0.15 \\
\bottomrule
\end{tabular}
\end{table}

% ---------------------------------------------------------------------
\subsection{Shock tubes with mechanical relaxation}
\label{sec:res_shock}
The following Riemann problems on $[0,1]$ (on $[-10,10]$ for the detonation products) with transmissive boundaries
are computed with the $p$ relaxation and compared with the exact solution of the five-equation model. The fluids are
at rest initially.

\paragraph{Water--air shock tube}
Water and air as in Sec.~\ref{sec:res_interface} (Saurel et al.\ 2009, Sec.~4.3.1). For $x<0.75$, water with
$\alpha_{\rm water}=1-10^{-6}$, $\rho=1000$, $p=10^9$\,Pa; for $x>0.75$, air with $\alpha_{\rm air}=1-10^{-6}$, $\rho=1$, $p=10^5$\,Pa;
$t_{\rm end}=240\,\mu$s, with 1000 and 100 cells. In the variant of Schropff et al.\ (2026) the interface is at $x=0.7$,
$\alpha_{\rm water}=1-10^{-7}$ or $10^{-7}$, $\rho_{\rm air}=50$ and $t_{\rm end}=241\,\mu$s, with 200, 1000 and 5000 cells.

\paragraph{Water--air shock tube at 1\,TPa}
The same fluids with $p=10^{12}$\,Pa in the water, $\rho_{\rm air}=10$ and the interface at $x=0.6$, the position consistent
with the wave positions reported by Saurel et al.\ (2009, Sec.~4.3.2); 1000 cells, $t_{\rm end}=8.3\,\mu$s.

\paragraph{Nitromethane shock tube}
Both phases are liquid nitromethane with the Cochran--Chan law (Saurel et al.\ 2009, Sec.~4.2),
\begin{equation}
p_r(\rho)=A_1x^{E_1}-A_2x^{E_2},\qquad
\varepsilon_r(\rho)=\frac1\rho\Big(\frac{A_1x^{E_1}}{E_1-1}-\frac{A_2x^{E_2}}{E_2-1}\Big),\qquad x=\frac\rho{\rho_0},
\label{eq:cochran}
\end{equation}
with $\Gamma=1.19$, $\rho_0=1134$, $A_1=0.819181$\,GPa, $A_2=1.50835$\,GPa, $E_1=4.52969$, $E_2=1.42144$. Phase 1 has
$\rho_1=1134$ and phase 2 $\rho_2=1200$; $\alpha_1=1-10^{-8}$, $p=2\times10^{10}$\,Pa for $x<0.5$ and $\alpha_1=10^{-8}$,
$p=2\times10^5$\,Pa otherwise. 500 cells, $t_{\rm end}=67\,\mu$s. The nonlinear reference curve makes the contact
discontinuity a test of pressure and velocity oscillations.

\paragraph{Detonation products--water shock tube}
Detonation products follow the JWL law (Pelanti 2022, Sec.~8.1.1),
\begin{equation}
p_r(\rho)=Ae^{-R_1\rho_0/\rho}+Be^{-R_2\rho_0/\rho},\qquad
\varepsilon_r(\rho)=\frac{A}{R_1\rho_0}e^{-R_1\rho_0/\rho}+\frac{B}{R_2\rho_0}e^{-R_2\rho_0/\rho},
\label{eq:jwl}
\end{equation}
with $\Gamma=0.25$, $\rho_0=1840$, $A=854.5$\,GPa, $B=20.5$\,GPa, $R_1=4.6$, $R_2=1.35$; water is NASG with $(\gamma,\pi,q,q',c_v,b)=
(1.187,\ 7.028\times10^8,\ -1177788,\ 0,\ 3610,\ 6.61\times10^{-4})$. On $[-10,10]$, the products ($\alpha=1-10^{-8}$, $\rho=2000$,
$p=4.6406\times10^{10}$\,Pa) fill $x<0$ and water ($\rho=1044$, $p=10^5$\,Pa) $x>0$. 500 cells, $t_{\rm end}=0.9$\,ms.

\paragraph{Epoxy--spinel mixture}
A mixture of epoxy (SG, $\gamma=2.43$, $\pi=5.3\times10^9$\,Pa) and spinel (SG, $\gamma=1.62$, $\pi=1.41\times10^{11}$\,Pa) with
$\alpha_{\rm epoxy}=0.5954$, $\rho_{\rm epoxy}=1185$, $\rho_{\rm spinel}=3622$ everywhere (Saurel et al.\ 2009, Secs.~5.1--5.2). The pressure
is $10^{10}$\,Pa for $x<0.6$ and $10^5$\,Pa otherwise, $t_{\rm end}=80\,\mu$s; in the second case it is $2\times10^{11}$\,Pa,
$t_{\rm end}=34\,\mu$s. 500 cells. The first case is also computed to $t_{\rm end}=30\,\mu$s with 200 and 1000 cells and $C=0.2$
(Schropff et al.\ 2026). Both phases are compressed through the same shock; the test checks the partition of
the energy defect of Sec.~\ref{sec:closure}. The first case uses $C=0.05$, the second the global Rusanov speed.

% ---------------------------------------------------------------------
\subsection{Expansion tubes and cavitation}
\label{sec:res_cav}
Two streams moving apart from $x=0.5$ in $[0,1]$, with transmissive boundaries, produce a strong expansion at the
centre. The same configuration is computed with the different relaxation levels and, for NASG water, with finite
rates.

\paragraph{Water--air expansion tube}
Water and air as in Sec.~\ref{sec:res_interface}, $\alpha_{\rm air}=10^{-2}$, $\rho_{\rm water}=1000$, $\rho_{\rm air}=1$, $p=10^5$\,Pa and
$u=-100$\,m/s for $x<0.5$, $+100$\,m/s otherwise (Saurel et al.\ 2009, Sec.~4.5); $p$ relaxation, 1000 cells,
$t_{\rm end}=1.85$\,ms. The exact solution of the five-equation model has a vacuum at the centre.

\paragraph{Water cavitation tube with phase transition}
Liquid water with vapour $\alpha_v=10^{-2}$, both phases at $p=10^5$\,Pa and the same temperature $T_0$, densities
$\rho_k=\rho_k(p,T_0)$, and $u=\mp U_0$ for $x\lessgtr0.5$.
\begin{itemize}
\item SG parameters (Pelanti and Shyue 2014, Sec.~5.1): liquid $(2.35,\ 10^9,\ -1167\times10^3,\ 0,\ 1816)$, vapour
$(1.43,\ 0,\ 2030\times10^3,\ -23.4\times10^3,\ 1040)$, $T_0=354.728$\,K. $U_0=2$\,m/s with $t_{\rm end}=3.2$\,ms, and $U_0=500$\,m/s
with $t_{\rm end}=0.58$\,ms; levels $p$, $p$--$pT$, $p$--$pT$--$pTg$, $p$--$pTg$.
\item NASG parameters (Pelanti 2022, Sec.~8.2.1): liquid $(1.187,\ 7.028\times10^8,\ -1177788,\ 0,\ 3610,\ 6.61\times10^{-4})$,
vapour $(1.467,\ 0,\ 2077616,\ 14317,\ 955,\ 0)$, $T_0=353$\,K, $U_0=2$\,m/s, $t_{\rm end}=3$\,ms; five levels: $p$; $p$--$pT(f)$ with
$\theta=2000$; $p$--$pT$; $p$--$pT$--$pTg(f)$ with $\nu=10^{-4}$; $p$--$pT$--$pTg$.
\end{itemize}
5000 cells in all cases.

% ---------------------------------------------------------------------
\subsection{Dodecane liquid--vapour shock tube}
\label{sec:res_dodecane}
Liquid dodecane at high pressure expands into its vapour (Pelanti and Shyue 2014, Sec.~5.2; Pelanti 2022, Sec.~8.2.2).
For $x<0.75$, liquid with $\alpha_l=1-10^{-8}$, $\rho_l=500$, $p=10^8$\,Pa; for $x>0.75$, vapour with $\alpha_v=1-10^{-8}$,
$\rho_v=2$, $p=10^5$\,Pa; transmissive boundaries, 2000 cells, $t_{\rm end}=473\,\mu$s. Two parameter sets are used.
\begin{itemize}
\item SG: liquid $(2.35,\ 4\times10^8,\ -775.269\times10^3,\ 0,\ 1077.7)$, vapour $(1.025,\ 0,\ -237.547\times10^3,\ -24.4\times10^3,\ 1956.45)$;
levels $p$, $p$--$pT$--$pTg$, $p$--$pTg$. $C=0.05$.
\item NASG: liquid $(1.206,\ 1.681\times10^8,\ -996054,\ 0,\ 2532,\ 7.51\times10^{-4})$, vapour $(1.021,\ 0,\ -384592,\ -4301,\ 2274,\ 0)$;
levels $p$ and $p$--$pT$--$pTg$.
\end{itemize}

% ---------------------------------------------------------------------
\subsection{Bart\'ak rapid depressurisation}
\label{sec:res_bartak}
A pipe of length 1.7\,m filled with hot pressurised water is opened at one end (Pelanti 2022, Sec.~8.2.4). Water is
NASG with liquid $(1.387,\ 8.899\times10^8,\ -1244191,\ 0,\ 3202,\ 4.78\times10^{-4})$ and vapour $(1.954,\ 0,\ 2287484,\ 6417,\ 462,\ 0)$.
Initially $p=12.5$\,MPa, $T=563.15$\,K, $u=0$, $\alpha_v=10^{-3}$ and $\rho_k=\rho_k(p,T)$. The right end $x=1.7$ is a wall; at the
left end the pressure $10^5$\,Pa is imposed: the ghost cells keep $\alpha$, $m_k$ and $u$ of the adjacent cell with both
phase pressures equal to $10^5$\,Pa, and this pressure is a given value in the pressure system. 1000 cells,
$t_{\rm end}=17.5$\,ms; pressure and vapour mass fraction are recorded at $x=48$\,mm. Three cases are computed with
instantaneous thermal relaxation: no mass transfer ($p$--$pT$), instantaneous mass transfer ($p$--$pT$--$pTg$), and the
finite rate
\begin{equation}
\nu=C_r\,\alpha_v^{0.6}\Big(\frac{p_{\rm sat}(T)-p}{p_{\rm crit}-p_{\rm sat}(T)}\Big)^{1.76}\quad(p<p_{\rm sat}(T)),\qquad \nu=0\ \text{otherwise},
\label{eq:bartak_nu}
\end{equation}
with $C_r=1$ and $p_{\rm crit}=22.064$\,MPa, evaluated in each cell from its $pT$-equilibrium state at the beginning of the
step and frozen over the stages. With instantaneous mass transfer the solution at the probe is a simple wave of the
homogeneous equilibrium model, which is computed exactly.

% ---------------------------------------------------------------------
\subsection{Shock--helium bubble interaction}
\label{sec:res_he}
A Mach 1.5 shock in air hits a helium bubble (Saurel et al.\ 2009, Sec.~4.6). The domain is $[0,0.3]\times[0,0.08]$ with
walls at $y=0$ and $y=0.08$ and transmissive boundaries in $x$; air is an ideal gas with $\gamma=1.4$, helium with
$\gamma=1.67$. The bubble has diameter 0.04 and centre $(0.2,0.04)$; inside $\alpha_{\rm air}=10^{-6}$, outside $1-10^{-6}$. Ahead
of the shock, air has $\rho=1.29$, $p=101325$\,Pa and helium $\rho=0.167$, both at rest. The shock is at $x=0.221$ and
moves to the left; behind it the air state follows from the Rankine--Hugoniot relations of the ideal gas at Mach 1.5.
$600\times160$ cells, $t_{\rm end}=600\,\mu$s, $p$ relaxation.

% ---------------------------------------------------------------------
\subsection{Cavitating Richtmyer--Meshkov instability}
\label{sec:res_rm}
A water column with a curved surface moving at $-200$\,m/s hits a wall (Saurel et al.\ 2009, Sec.~4.7). On $[0,3]\times[0,1]$,
water (SG, $\gamma=4.4$, $\pi=6\times10^8$\,Pa, $\rho=1000$) occupies $x<1.2+\sqrt{0.36-(y-0.5)^2}$ and gas (ideal, $\gamma=1.8$,
$\rho=100$) the rest, with $\alpha=1-10^{-6}$ for the present phase; everywhere $u=-200$\,m/s and $p=101325$\,Pa. The left,
bottom and top boundaries are walls and the right boundary is transmissive. $300\times100$ cells, $t_{\rm end}=8.6$\,ms,
$p$ relaxation, global Rusanov speed.

% ---------------------------------------------------------------------
\subsection{Compression of a vapour bubble}
\label{sec:res_bubble}
A dodecane vapour bubble in liquid dodecane is compressed by a piston (Pelanti and Shyue 2014, Sec.~5.4). The SG
parameters are those of Sec.~\ref{sec:res_dodecane}. On $[0,1]^2$ a bubble of radius 0.2 is centred at $(0.5,0.5)$;
initially $p=10^5$\,Pa and $T=600$\,K, so $\rho_l=458.338$ and $\rho_v=3.408$, with $\alpha_l=10^{-6}$ inside the bubble and
$\alpha_v=10^{-6}$ outside, at rest. All boundaries are walls; the left wall moves with velocity $100$\,m/s in $+x$ (the grid
is fixed, and the wall boundary condition imposes this normal velocity). $200\times200$ cells, $t_{\rm end}=1.2$\,ms,
$p$--$pT$--$pTg$ relaxation, global Rusanov speed.

% ---------------------------------------------------------------------
\subsection{Underwater explosion near a wall}
\label{sec:res_undex}
A bubble of detonation products expands in water below a rigid wall (Pelanti 2022, Sec.~8.1.2). The states and the
equations of state are those of the detonation products--water shock tube of Sec.~\ref{sec:res_shock}: products with
$\rho=2000$, $p=4.6406\times10^{10}$\,Pa and $\alpha=1-10^{-6}$ inside a circle of radius 0.05 centred at $(0,-0.32)$, water with
$\rho=1044$, $p=10^5$\,Pa outside. The domain is $[-1.35,1.35]\times[-1.5,0]$ with a wall at $y=0$ and transmissive
boundaries elsewhere. $361\times200$ cells, $t_{\rm end}=0.4$\,ms, $p$ relaxation; the pressure is recorded at the wall
centre $(0,0)$.
